\documentclass[../../main.tex]{subfiles}
\begin{document}

{\bf [解答]}

題意より$a = 0.1\alpha$，$b = 0.1\beta$ とおく（$\alpha, \beta \in \mathbb{N}$）．距離 $x\,\mathrm{km}$ のときの運賃を $f(x)$ 円とすると，
\begin{align}
  \begin{cases}
    f(319) = 970 \\
    f(349) = 1010
  \end{cases} \label{eq:1}
\end{align}
である．題意の運賃計算規則から，
\begin{itemize}
  \item $f(319)$ は $300a + 19b = 30\alpha + 1.9\beta$ の10円未満の端数を切り上げたもの
  \item $f(349)$ は $300a + 49b = 30\alpha + 4.9\beta$ の10円未満の端数を切り上げたもの
\end{itemize}
であるから，
\begin{align}
  \begin{cases}
    960 < 30\alpha + 1.9\beta \le 970 \\
    1000 < 30\alpha + 4.9\beta \le 1010
  \end{cases} \label{eq:2}
\end{align}
が成り立つ．\cref{eq:2}の辺々を引くと，
\begin{align}
  30 < 3\beta < 50 \quad \therefore 10 < \beta < \dfrac{50}{3} = 16.6\dots
\end{align}
となる．$\beta \in \mathbb{N}$ より，$\beta$ の候補は
\begin{align}
  \beta = 11, 12, 13, 14, 15, 16
\end{align}
である．ここで実運賃の差を $A = f(349) - f(319)$ とおくと，\cref{eq:1}より $A = 1010 - 970 = 40$ となる．各 $\beta$ がこの条件を満たすか調べる．

\subparagraph{$1^\circ$ $\beta = 11$ のとき}
$1.9\beta = 20.9$ より切り上げは $30$ 円，$4.9\beta = 53.9$ より切り上げは $60$ 円となるから，
\begin{align}
  A = 60 - 30 = 30 \ne 40
\end{align}
となり不適．

\subparagraph{$2^\circ$ $\beta = 12$ のとき}

$1.9\beta = 22.8$ より切り上げは $30$ 円，$4.9\beta = 58.8$ より切り上げは $60$ 円となるから，
\begin{align}
  A = 60 - 30 = 30 \ne 40
\end{align}
となり不適．

\subparagraph{$3^\circ$ $\beta = 13$ のとき}

$1.9\beta = 24.7$ より切り上げは $30$ 円，$4.9\beta = 63.7$ より切り上げは $70$ 円となり $A = 70 - 30 = 40$ を満たす．しかし，\cref{eq:1}より
\begin{align}
  970 = 30\alpha + 30 \quad \therefore \alpha = \dfrac{94}{3} \notin \mathbb{N}
\end{align}
となり不適．

\subparagraph{$4^\circ$ $\beta = 14$ のとき}

$1.9\beta = 26.6$ より切り上げは $30$ 円，$4.9\beta = 68.6$ より切り上げは $70$ 円となり $A = 70 - 30 = 40$ を満たすが，$3^\circ$ と同様に\cref{eq:1}より
\begin{align}
  970 = 30\alpha + 30 \quad \therefore \alpha = \dfrac{94}{3} \notin \mathbb{N}
\end{align}
となり不適．

\subparagraph{$5^\circ$ $\beta = 15$ のとき}

$1.9\beta = 28.5$ より切り上げは $30$ 円，$4.9\beta = 73.5$ より切り上げは $80$ 円となるから，
\begin{align}
  A = 80 - 30 = 50 \ne 40
\end{align}
となり不適．

\subparagraph{$6^\circ$ $\beta = 16$ のとき}
$1.9\beta = 30.4$ より切り上げは $40$ 円，$4.9\beta = 78.4$ より切り上げは $80$ 円となり，$A = 80 - 40 = 40$ を満たす．\cref{eq:1}に代入すると，
\begin{align}
  970 = 30\alpha + 40 \quad \therefore \alpha = 31 \in \mathbb{N}
\end{align}
を得る．このとき $f(349) = 30 \times 31 + 80 = 1010$ も満たし，適する．
\casesend

以上より，$(\alpha, \beta) = (31, 16)$，すなわち
\begin{align}
  a = 3.1, \quad b = 1.6 \tag{答}
\end{align}
である．

\newpage
\end{document}
